\Article{4}
Neben den magnetischen Kräften, die einzelne Körper, wie wir sehen, bei hinreichender Nähe aufeinander ausüben, wirkt eine weitere Kraft auf die magnetischen Flüssigkeiten. Da sie sich überall auf der Erde zeigt, schreiben wir sie der Erdkugel selbst zu und nennen sie Erdmagnetismus. Diese Kraft äußert sich auf zweierlei Weise. Werden Körper der zweiten Art, in denen der Magnetismus entwickelt ist, in ihrem Schwerpunkt gestützt, so sucht die Kraft sie in eine bestimmte Richtung zu drehen. In Körpern der ersten Art dagegen trennt sie die magnetischen Flüssigkeiten von selbst. Wählt man Körper von geeigneter Gestalt und bringt sie in eine geeignete Lage, so lässt sich diese Trennung sehr deutlich wahrnehmen. Beide Erscheinungen erklären sich, wenn man annimmt, dass diese Kraft an jedem Ort die nördliche magnetische Flüssigkeit in eine bestimmte Richtung und die südliche mit gleicher Intensität in die entgegengesetzte Richtung treibt. Mit der Richtung des Erdmagnetismus meinen wir stets die erstere. Sie ist demnach durch ihre Neigung zur \SourceLeaf{93}Horizontalebene und durch die Abweichung der Vertikalebene, in der sie wirkt, von der Meridianebene bestimmt. Jene Vertikalebene heißt \textit{magnetische Meridianebene}. Die Intensität des Erdmagnetismus ist an der bewegenden Kraft zu messen, die er auf eine Einheit der freien magnetischen Flüssigkeit ausübt.

Diese Kraft ist an verschiedenen Orten der Erde verschieden und ändert sich auch am selben Ort im Lauf der Jahrhunderte und Jahre. Sie ändert sich ferner im Lauf der Sommer und der Tagesstunden. Für die Richtung ist diese Veränderlichkeit seit Langem bekannt. Für die Intensität ließ sie sich bisher nur im Lauf der Tagesstunden beobachten, weil uns für längere Zeiträume geeignete Hilfsmittel fehlten. Die Angabe der Intensität in absoluten Einheiten wird diesem Mangel künftig abhelfen.

\Article{5}
Um die Wirkung des Erdmagnetismus auf magnetische Körper der zweiten Art zu berechnen, die wir von nun an stets voraussetzen, denke man sich einen solchen Körper in unendlich kleine Teile zerlegt. Es sei $\mathrm{d}m$ das Element des freien Magnetismus in einem Teilchen, dessen Koordinaten bezüglich dreier zueinander senkrechter und gegenüber dem Körper fester Ebenen wir mit $x,y,z$ bezeichnen. Die Elemente der südlichen Flüssigkeit rechnen wir negativ. Zunächst ist also klar, dass das über den ganzen Körper, ja über jeden messbaren Teil des Körpers erstreckte Integral $\int\mathrm{d}m$ den Wert $=0$ hat. Setzen wir $\int x\mathrm{d}m=X$, $\int y\mathrm{d}m=Y$, $\int z\mathrm{d}m=Z$. Diese Größen können wir die Momente des freien Magnetismus bezüglich der drei Grundebenen oder bezüglich der darauf senkrechten Achsen nennen. Bezeichnet $a$ eine beliebige Konstante, so gilt $\int(x-a)\mathrm{d}m=X$. Das Moment bezüglich einer gegebenen Achse hängt also nur von ihrer Richtung ab und ist von ihrem Ursprung unabhängig. Ziehen wir durch den Koordinatenursprung eine vierte Achse, die mit den ursprünglichen Achsen die Winkel $A,B,C$ bildet, so ist das Moment des Elements $\mathrm{d}m$ bezüglich dieser Achse $=(x\cos A+y\cos B+z\cos C)\mathrm{d}m$ und damit das Moment des freien Magnetismus im ganzen Körper
\[
=X\cos A+Y\cos B+Z\cos C=V
\]
Man setze
\[
\BareRoot{(XX+YY+ZZ)}=M,\enspace\text{und }X=M\cos\alpha,\enspace Y=M\cos\beta,\quad Z=M\cos\gamma
\]
und ziehe eine fünfte Achse, die mit den drei ursprünglichen die Winkel $\alpha,\beta,\gamma$ und mit der vierten den Winkel $\omega$ bildet. Da bekanntlich gilt
\SourceLeaf{94}
\[
\cos\omega=\cos A\cos\alpha+\cos B\cos\beta+\cos C\cos\gamma,\quad\text{ergibt sich }V=M\cos\omega
\]
Wir nennen diese fünfte Achse einfach die \textit{magnetische Achse} des Körpers und nehmen an, dass ihre \textit{Richtung} dem positiven Wert der Wurzel $\BareRoot{(XX+YY+ZZ)}$ entspricht. Fällt die vierte Achse mit dieser magnetischen Achse zusammen, so wird das Moment $V$ zu $=M$, dem offenbar größten aller Momente. Das Moment bezüglich jeder anderen Achse findet man, indem man dieses größte Moment mit dem Kosinus des Winkels zwischen jener Achse und der magnetischen Achse multipliziert. Wo keine Verwechslung zu befürchten ist, können wir das größte Moment einfach das Moment des Magnetismus nennen. Das Moment bezüglich jeder zur magnetischen Achse senkrechten Achse ist $=0$. Bezüglich jeder Achse, die mit der magnetischen Achse einen stumpfen Winkel bildet, ist es negativ.

Die magnetische Achse ist also keine festgelegte Gerade. Da sie durch jeden beliebigen Punkt gezogen werden kann, legt sie nur eine Richtung fest. Es gibt mithin unendlich viele zueinander parallele magnetische Achsen. Wählen wir eine davon nach Belieben und geben ihr eine bestimmte Länge, so heißen ihre Enden Pole. Die Achse ist vom Südpol zum Nordpol gerichtet.

\Article{6}
Wirkt auf die einzelnen Teilchen der magnetischen Flüssigkeiten eine Kraft von konstanter Intensität und Richtung, so lässt sich die daraus entstehende Gesamtkraft auf den Körper leicht aus den Grundsätzen der Statik ableiten. Bei den hier betrachteten Körpern haben diese Teilchen ihre Fließfähigkeit gewissermaßen verloren und bilden mit dem wägbaren Körper eine einzige starre Masse. Auf ein beliebiges magnetisches Molekül $\mathrm{d}m$ wirke eine bewegende Kraft $=P\mathrm{d}m$ in Richtung $D$. Bei den Molekülen der südlichen Flüssigkeit bedeutet das negative Vorzeichen schon für sich die entgegengesetzte Richtung. Es seien $A,B$ zwei Punkte des Körpers, die in Richtung der magnetischen Achse liegen, und ihr Abstand sei $=r$, positiv gerechnet, wenn die magnetische Achse von $A$ nach $B$ verläuft. Fügt man zu diesen Kräften zwei neue hinzu, jede $=\frac{PM}{r}$, von denen die eine bei $A$ in Richtung $D$ und die andere bei $B$ in der entgegengesetzten Richtung wirkt, so sind leicht ersichtlich alle diese Kräfte im Gleichgewicht. Die ursprünglichen Kräfte sind daher zwei Kräften $=\frac{PM}{r}$ gleichwertig, von denen die eine bei $B$ in Richtung $D$ und die andere bei $A$ in der entgegengesetzten Richtung wirkt. Diese beiden Kräfte lassen sich offenbar nicht zu einer einzigen zusammenfassen.

Wirkt außer der Kraft $P$ eine weitere gleichartige Kraft $P'$ in Richtung $D'$ auf die magnetischen Flüssigkeiten des Körpers, so lassen sich an ihrer Stelle wiederum zwei andere einsetzen. Sie können entweder bei denselben Punkten $A,B$ oder \SourceLeaf{95}allgemeiner bei zwei anderen Punkten $A',B'$ wirken, sofern $A'B'$ ebenfalls eine magnetische Achse ist. Setzt man den Abstand $A'B'=r'$, so müssen diese Kräfte $=\frac{P'M}{r'}$ sein. Bei $B'$ wirkt die Kraft in Richtung $D'$, bei $A'$ in der entgegengesetzten Richtung. Entsprechendes gilt für mehrere Kräfte.

Innerhalb eines so kleinen Raums, wie ihn der für den Versuch bestimmte Körper einnimmt, darf man die Intensität und Richtung der erdmagnetischen Kraft überall als konstant ansehen, auch wenn sie sich mit der Zeit ändern können. Daher lässt sich das eben Gesagte auf diese Kraft anwenden. Es kann jedoch zweckmäßig sein, sie gleich von Anfang an in zwei Kräfte zu zerlegen, eine horizontale $=T$ und eine vertikale $=T'$, die in unseren Gegenden nach unten gerichtet ist. Will man für die letztere zwei andere Kräfte einsetzen, die bei $A',B'$ wirken, so kann man sowohl den Punkt $A'$ als auch den Abstand $A'B'=r'$ beliebig wählen. Wir wählen den Schwerpunkt als $A'$ und bezeichnen das Gewicht des Körpers, also die bewegende Kraft, die die Schwerkraft auf seine Masse ausübt, mit $p$. Dann setzen wir $\frac{T'M}{p}=r'$. So zerfällt die Wirkung der Kraft $T'$ in eine nach oben gerichtete Kraft $=p$ bei $A'$ und eine gleich große, nach unten gerichtete Kraft bei $B'$. Da die Schwerkraft selbst die erstere offenbar aufhebt, läuft die Wirkung der vertikalen erdmagnetischen Kraft einfach darauf hinaus, dass der Schwerpunkt von $A'$ nach $B'$ verlegt wird. In den Gegenden, in denen die erdmagnetische Kraft einen spitzen Winkel mit der Vertikalen bildet, also ihr vertikaler Anteil die nördliche magnetische Flüssigkeit nach oben treibt, findet offenbar eine entsprechende Verlegung des Schwerpunkts auf der magnetischen Achse in Richtung des Südpols statt.

Aus dieser Auffassung ergibt sich unmittelbar, dass man die Inklination aus Versuchen mit einer Magnetnadel in einem \textit{einzigen} magnetischen Zustand allein nicht ableiten kann, gleichgültig, welche Versuche man anstellt. Vielmehr muss die Lage des \textit{wirklichen} Schwerpunkts bereits anderweitig bekannt sein. Gewöhnlich bestimmt man diese Lage, bevor man die Nadel magnetisiert. Dieses Vorgehen ist jedoch wenig zuverlässig, da eine Stahlnadel meist schon bei ihrer Herstellung einen wenn auch schwachen Magnetismus annimmt. Um die Inklination zu bestimmen, muss man daher durch eine geeignete Änderung des magnetischen Zustands der Nadel eine \textit{andere} Verlegung des Schwerpunkts herbeiführen. Damit sie sich möglichst stark von der ersten unterscheidet, muss man die Pole umkehren. Auf diese Weise lässt sich eine doppelte Verschiebung erzielen. Selbst bei Nadeln mit den günstigsten Abmessungen, die bis zur Sättigung magnetisiert sind, kann die Verlegung des Schwerpunkts eine bestimmte Grenze nicht überschreiten. Für eine einfache Verschiebung beträgt sie in unseren Gegenden etwa $0{,}4$ Millimeter und \SourceLeaf{96}bleibt unter $0{,}6$ Millimetern in Gegenden, in denen die vertikale Kraft am größten ist. Daran erkennt man zugleich, wie viel mechanische Präzision bei Nadeln erforderlich ist, die zur Bestimmung der Inklination dienen sollen.

\Article{7}
Ein Punkt $C$ eines magnetischen Körpers sei fest. Für das Gleichgewicht ist notwendig und hinreichend, dass die durch $C$, den Schwerpunkt und die magnetische Achse gelegte Ebene mit der magnetischen Meridianebene zusammenfällt und dass sich außerdem die Momente aufheben, mit denen die erdmagnetische Kraft und die Schwerkraft jene Ebene um $C$ zu drehen suchen. Die zweite Bedingung läuft auf Folgendes hinaus. Bezeichnet $T$ den horizontalen Anteil der erdmagnetischen Kraft und $i$ die Neigung der magnetischen Achse zur Horizontalebene, so muss $TM\sin i$ gleich dem Produkt aus dem Gewicht des Körpers und dem Abstand des verlegten Schwerpunkts $B'$ von der durch $C$ gezogenen Vertikalen sein. Offenbar muss dieser Abstand auf der südlichen beziehungsweise nördlichen Seite liegen, je nachdem, ob $i$ eine Erhebung oder eine Senkung ist. Für $i=0$ liegt $B'$ auf der Vertikalen selbst. Der Körper werde nun um diese Vertikale so gedreht, dass die magnetische Achse in eine Vertikalebene gelangt, deren magnetisches Azimut, also deren Winkel mit dem nördlichen Teil des magnetischen Meridians, $=u$ ist. Dabei darf man Winkel nach Osten oder nach Westen nach Belieben als positiv rechnen. Der Erdmagnetismus übt dann eine Kraft aus, die den Körper um die vertikale Achse zu drehen, also den Winkel $u$ zu verkleinern sucht. Ihr Moment ist $=TM\cos i\sin u$, und der Körper führt um diese Achse Schwingungen aus, deren Dauer sich nach bekannten Methoden berechnen lässt. Bezeichnen wir nämlich mit $K$ das Trägheitsmoment des Körpers bezüglich der Schwingungsachse, also die Summe der wägbaren Moleküle, jeweils mit dem Quadrat ihres Abstands von der Achse multipliziert, und wie üblich mit $\pi$ den halben Umfang eines Kreises vom Radius $=1$, so beträgt die Dauer einer unendlich kleinen Schwingung $=\pi\BareRoot{\frac{K}{TM\cos i}}$. Dabei muss den Größen $T,M$ diejenige Einheit der beschleunigenden Kräfte zugrunde liegen, die in einer Zeiteinheit die Geschwindigkeit $=1$ erzeugt. Die Zurückführung endlicher Schwingungen auf unendlich kleine lässt sich ähnlich berechnen wie bei den Schwingungen eines Pendels. Ergibt sich also aus den Beobachtungen für die Dauer einer unendlich kleinen Schwingung der Wert $=t$, so haben wir $TM=\frac{\pi\pi K}{tt\cos i}$. Ist der Körper nun so aufgehängt, dass die magnetische Achse horizontal liegt, wie wir von nun an stets voraussetzen, so gilt
\[
TM=\frac{\pi\pi K}{tt}
\]
Möchte man lieber die Schwerkraft als Einheit der beschleunigenden Kräfte wählen, so müsste man diesen \SourceLeaf{97}Wert durch $\pi\pi l$ teilen. Dabei bezeichnet $l$ die Länge eines einfachen Pendels, das in einer Zeiteinheit eine Schwingung ausführt. Allgemein ergibt sich damit $TM=\frac{K}{ttl\cos i}$, für unseren Fall also $TM=\frac{K}{ttl}$.
